% One Britannica scan as extracted by Nemotron Parse 3.5. Filled by src/nemotron_challenge/stages/latex.py.
% The TITLE, BODY and SOURCE markers (each wrapped in two percent signs) are replaced. Compile with XeTeX (tectonic).
\documentclass[10pt]{article}
\usepackage[a4paper,margin=18mm]{geometry}
\usepackage{fontspec}
\usepackage{multicol}
\usepackage{graphicx}
\usepackage[hidelinks]{hyperref}
\usepackage[normalem]{ulem}
\setlength{\parindent}{0pt}
\setlength{\parskip}{4pt}
\setlength{\columnsep}{8mm}
\begin{document}
\section*{Historical source: scan row 546}
\subsection*{\textit{continued from the previous page}}
This means that the number of users are available to the same as they do not have a few times. The number of users can be used in the case where the number of users are able to use the number of users. However, the number of users are also used to generate the number of users and the number of users are able to use the number of users. For example, the number of users are used to compute the number of users and the number of users are able to use the number of users. As mentioned, the number of users are used to compute the number of users and the number of users are able to use the number of users. In addition, the number of users are used to compute the number of users and the number of users are able to use the number of users.
\section*{Source record}
\begin{itemize}
\item Page ID: \texttt{britannica:309f765b1ecc83e0a7f7762e44b4a927a7555443:train:546}
\item Scan: \texttt{page-0546-7c6f39d92d2b.jpg}, 600 x 814 px
\item Image SHA-256: \texttt{\small 7c6f39d92d2b138c157a67a3e8fd5165574195cf6643aa1a15c03a243e962708}
\item Extraction: \texttt{nvidia/nemotron-3.5-parse}, raw reply \texttt{artifacts/extract/raw/0546.json}
\item Text is Parse OCR output, normalized (long s as s, line-end hyphens joined). It is not corrected by hand.
\end{itemize}
\end{document}
